% Rozdzial 19. Wymaga makr \code, \term, \nota oraz pakietow:
% listings, tabularx, booktabs, enumitem, hyperref, amsmath.
% Fragment rozdzialu (bez preambuly). C++20; nowsze API zaznaczono.
\chapter{Wejście, wyjście, system plików i przetwarzanie tekstu}
\label{chap:cpp19-io}

\section{Wprowadzenie: dane, reprezentacja i granice odpowiedzialności}

Operacje wejścia i wyjścia (\emph{I/O}) przenoszą dane między programem
oraz jego otoczeniem: terminalem, plikiem, potokiem czy urządzeniem.
W odróżnieniu od zwykłych obliczeń mogą zakończyć się częściowym sukcesem,
błędem systemu albo blokowaniem. Dobrze zaprojektowany kod oddziela trzy
warstwy: \emph{pozyskanie bajtów}, \emph{parsowanie reprezentacji} oraz
\emph{przetworzenie obiektów}. Dzięki temu test parsera nie wymaga dysku,
a zmiana formatu pliku nie wymusza zmiany algorytmu numerycznego.

\begin{center}\small
\begin{tabularx}{\textwidth}{@{}p{0.28\textwidth}X@{}}
\toprule
Mechanizm & Typowe zastosowanie \\
\midrule
\code{iostream}, \code{fstream} & Formatowany tekst, prosta obsługa plików, strumienie. \\
\code{filesystem} & Ścieżki, katalogi, metadane, operacje na plikach. \\
\code{string}, \code{string_view} & Własność tekstu i nieposiadający go widok. \\
\code{charconv} & Szybkie parsowanie i formatowanie liczb bez ustawień regionalnych. \\
\code{format} (C++20) & Składanie tekstu według opisu formatowania. \\
\code{print} (C++23) & Wygodne formatowane wypisywanie; zależy od implementacji biblioteki. \\
\bottomrule
\end{tabularx}
\end{center}

\nota{Przykłady są pisane dla C++20. W tym rozdziale \code{std::filesystem}
nie oznacza, że każda operacja jest atomowa ani że istnieją uprawnienia
do każdego wskazanego pliku. Kod obsługujący niezaufane ścieżki wymaga
osobnej analizy bezpieczeństwa.}

\section{Strumienie: stan, formatowanie i buforowanie}
\label{sec:cpp19-streams}

\code{std::istream} jest interfejsem odczytu, a \code{std::ostream}
interfejsem zapisu. \code{std::cin}, \code{std::cout}, \code{std::cerr}
i \code{std::clog} to gotowe strumienie standardowe. Operatory
\code{>>} i \code{<<} wykonują zwykle operacje \emph{formatowane}.
Strumień przechowuje także bity stanu: \code{goodbit}, \code{eofbit},
\code{failbit} i \code{badbit}. Stan jest obserwowalny przez
\code{good()}, \code{eof()}, \code{fail()}, \code{bad()} i \code{rdstate()}.

\subsection{Poprawna pętla odczytu}

Nie należy pisać \code{while (!in.eof())}: wskaźnik końca pliku jest
ustawiany \emph{po nieudanej próbie odczytu}, a nie przed nią.
Warunkiem pętli powinien być sam odczyt.

\begin{lstlisting}[caption={Odczyt liczb i kontrola bledu formatu (kompletny program, C++20).}]
#include <iostream>
#include <sstream>

int main() {
    std::istringstream in("10 20 30");
    int value = 0;
    int sum = 0;
    while (in >> value)
        sum += value;

    if (!in.eof()) {
        std::cerr << "Invalid input\n";
        return 1;
    }
    std::cout << sum << '\n'; // 60
}
\end{lstlisting}

Odczyt \code{int} ignoruje początkowe białe znaki. Gdy napotka
\code{abc}, ustawia \code{failbit}; naprawienie \emph{stanu strumienia}
przez \code{clear()} nie usuwa automatycznie błędnego tekstu z bufora.
Połączenie \code{>>} z \code{getline} wymaga uwzględnienia pozostawionego
znaku końca linii.

\begin{lstlisting}[caption={Liczba, a potem pelna linia (kompletny program).}]
#include <iostream>
#include <limits>
#include <sstream>
#include <string>

int main() {
    std::istringstream in("7\nhello world\n");
    int id{};
    in >> id;
    in.ignore(std::numeric_limits<std::streamsize>::max(), '\n');
    std::string line;
    std::getline(in, line);
    std::cout << id << ":" << line << '\n';
}
\end{lstlisting}

\subsection{Flagi formatowania i wyjątki strumieni}

Manipulatory \code{std::fixed}, \code{std::scientific},
\code{std::setprecision}, \code{std::setw}, \code{std::setfill} i
\code{std::boolalpha} zmieniają sposób prezentacji. Część ustawień
pozostaje w strumieniu na kolejne operacje; \code{setw} dotyczy tylko
następnego pola. Ustawienie \code{exceptions(mask)} sprawia, że zmiana
odpowiednich bitów stanu może zgłosić \code{std::ios_base::failure}.

\begin{lstlisting}[caption={Lokalne formatowanie bez zmieniania std::cout.}]
#include <iomanip>
#include <iostream>
#include <sstream>

int main() {
    std::ostringstream out;
    out << std::fixed << std::setprecision(2) << 3.14159;
    std::cout << out.str() << '\n'; // 3.14
}
\end{lstlisting}

\nota{\code{std::endl} nie jest tylko nową linią: dodatkowo opróżnia
bufor. W pętli wypisującej wiele rekordów używamy zwykle \code{'\\n'},
a \code{flush} wywołujemy, gdy rzeczywiście wymagamy wypchnięcia bufora.
\code{flush} nie oznacza trwałego zapisania danych na fizycznym nośniku.}

\subsection{Formatowanie a ustawienia regionalne}

Operatory strumieni zależą od przypisanej im lokalizacji
(\code{std::locale}); separator dziesiętny i grupowanie cyfr mogą się
różnić. Do stabilnego formatu wymiany danych należy jawnie wybrać
konwencję, najlepiej niezależną od regionalnych ustawień użytkownika.
\code{std::format} (C++20) zapewnia formatowanie według tekstu formatu,
lecz dostępność nagłówka \code{<format>} zależy od biblioteki.

\begin{lstlisting}[caption={Formatowanie C++20 -- wymaga obslugi <format>.}]
#include <format>
#include <string>

std::string describe(int n, double x) {
    return std::format("n={}, x={:.3f}", n, x);
}
\end{lstlisting}

\section{Pliki tekstowe i binarne}
\label{sec:cpp19-files}

\code{std::ifstream} czyta plik, \code{std::ofstream} zapisuje,
a \code{std::fstream} obsługuje oba kierunki. Otwarcie odbywa się
w konstruktorze lub przez \code{open}; destruktor zamyka uchwyt
zgodnie z RAII. Tryby obejmują \code{std::ios::in}, \code{out},
\code{app}, \code{trunc}, \code{binary} i \code{ate}.

\begin{lstlisting}[caption={Zapis i odczyt pliku tekstowego (kompletny program).}]
#include <fstream>
#include <iostream>
#include <stdexcept>
#include <string>

int main() {
    const std::string path = "cpp19_example.txt";
    {
        std::ofstream out(path, std::ios::trunc);
        if (!out) throw std::runtime_error("open for write");
        out << "alpha\nbeta\n";
        out.close();
        if (!out) throw std::runtime_error("write/close failed");
    }
    std::ifstream in(path);
    if (!in) throw std::runtime_error("open for read");
    std::string line;
    while (std::getline(in, line))
        std::cout << line << '\n';
    if (in.bad()) throw std::runtime_error("read failed");
}
\end{lstlisting}

Sprawdzanie sukcesu otwarcia nie wystarcza: zapis może zakończyć się
błędem także po częściowym wykonaniu. Jawne \code{close()} umożliwia
sprawdzenie stanu przed destrukcją. Sam fakt \code{close()} nie daje
gwarancji trwałości danych na dysku (\emph{durability}).

\subsection{Binarna reprezentacja obiektu nie jest formatem pliku}

\code{std::ios::binary} usuwa pewne transformacje trybu tekstowego;
nie wprowadza przenośnego sposobu kodowania struktur. Zapis
\code{reinterpret_cast<const char*>(&obj)} pomija kwestie
\emph{endianness}, wyrównania, paddingu, rozmiaru typów i wersji
reprezentacji. W plikach wymiany zalecamy własny jawny format albo
przetestowaną bibliotekę serializacji.

\begin{lstlisting}[caption={Jawny zapis liczby 32-bitowej w kolejnosci little-endian (C++20).}]
#include <array>
#include <cstdint>
#include <fstream>
#include <stdexcept>

void write_u32_le(std::ostream& out, std::uint32_t v) {
    std::array<char, 4> bytes{};
    for (unsigned i = 0; i < 4; ++i)
        bytes[i] = static_cast<char>((v >> (8 * i)) & 0xffu);
    out.write(bytes.data(), bytes.size());
    if (!out) throw std::runtime_error("write failed");
}
\end{lstlisting}

Odbiorca powinien sprawdzać długość danych, wersję protokołu,
dozwolone zakresy i limity alokacji, zanim zacznie tworzyć obiekty.
\code{std::bit_cast} ułatwia analizę reprezentacji typów o odpowiednich
własnościach, ale nie rozwiązuje samoistnie problemu przenośnej serializacji.

\section{System plików: \texttt{std::filesystem}}
\label{sec:cpp19-filesystem}

Nagłówek \code{<filesystem>} (C++17) udostępnia typ \code{path}
oraz operacje na ścieżkach i katalogach. \code{path} reprezentuje
ścieżkę w notacji systemu, nie otwiera jeszcze pliku. Operator
\code{/} poprawnie dołącza elementy ścieżki.

\begin{lstlisting}[caption={Sciezki i katalogi (kompletny program).}]
#include <filesystem>
#include <iostream>

namespace fs = std::filesystem;

int main() {
    const fs::path base = fs::temp_directory_path();
    const fs::path path = base / "results" / "data.csv";
    std::cout << path.filename().string() << '\n';
    std::cout << path.extension().string() << '\n';
    std::cout << path.parent_path().filename().string() << '\n';
}
\end{lstlisting}

\subsection{Wyjątki i \texttt{std::error\_code}}

Wiele funkcji ma przeciążenia zgłaszające
\code{std::filesystem::filesystem_error} oraz warianty przyjmujące
\code{std::error_code&}. Drugie przydają się w kodzie, w którym brak
pliku jest normalnym wynikiem, a nie wyjątkiem.

\begin{lstlisting}[caption={Kontrolowany odczyt informacji o pliku.}]
#include <filesystem>
#include <iostream>
#include <system_error>

int main() {
    std::error_code ec;
    const bool exists = std::filesystem::exists("config.txt", ec);
    if (ec) {
        std::cerr << "filesystem error: " << ec.message() << '\n';
        return 1;
    }
    std::cout << std::boolalpha << exists << '\n';
}
\end{lstlisting}

\subsection{Iterowanie i symbole łączy}

\code{directory_iterator} przechodzi po jednym poziomie, natomiast
\code{recursive_directory_iterator} schodzi do podkatalogów.
\code{exists}, \code{is_regular_file}, \code{is_symlink}, \code{status},
\code{symlink_status}, \code{canonical} i \code{weakly_canonical}
mają różne znaczenie. Domyślnie nie należy utożsamiać ścieżki
tekstowej z tożsamością rzeczywistego pliku: łącza symboliczne,
uprawnienia i zmiany katalogów mogą zmienić wynik.

\begin{lstlisting}[caption={Wypisanie zwyklych plikow z katalogu (C++17).}]
#include <filesystem>
#include <iostream>
#include <system_error>

int main() {
    namespace fs = std::filesystem;
    std::error_code ec;
    for (fs::directory_iterator it(fs::current_path(), ec), end;
         !ec && it != end; it.increment(ec)) {
        std::error_code item_ec;
        if (it->is_regular_file(item_ec) && !item_ec)
            std::cout << it->path().filename().string() << '\n';
    }
    if (ec) {
        std::cerr << ec.message() << '\n';
        return 1;
    }
}
\end{lstlisting}

\nota{Wzorzec \emph{sprawdź-then-użyj} nie zapobiega wyścigom
\emph{TOCTOU}: między \code{exists(path)} i \code{open(path)} ktoś może
zmienić obiekt wskazywany przez ścieżkę. Przy wymaganiach
bezpieczeństwa potrzebne są mechanizmy systemowe pracujące na
uchwytach do otwartych plików i właściwe polityki dostępu.}

\subsection{Zapis przez plik tymczasowy i podmianę}

Częsty wzorzec to zapis nowej zawartości do osobnego pliku w tym
samym katalogu, zamknięcie go i zastąpienie pliku docelowego przez
\code{rename}. Dokładne gwarancje atomowości podmiany, obsługę
istniejącego celu oraz trwałości trzeba ustalić dla danego systemu
plików i systemu operacyjnego. Sama biblioteka C++ nie tworzy
transakcyjnego systemu plików. Dla plików poufnych konieczne jest
również bezpieczne tworzenie unikalnego pliku tymczasowego.

\section{Własność tekstu i \texttt{std::string\_view}}
\label{sec:cpp19-text}

\code{std::string} zarządza własną pamięcią. \code{std::string_view}
przechowuje wskaźnik i długość fragmentu znaków, nie przedłużając
czasu życia danych. Widok bywa idealnym parametrem funkcji, jeżeli
funkcja nie zachowuje go na później.

\begin{lstlisting}[caption={Widok do bufora -- bez kopiowania (kompletny program).}]
#include <iostream>
#include <string>
#include <string_view>

void print_prefix(std::string_view text, std::size_t n) {
    std::cout << text.substr(0, n) << '\n';
}

int main() {
    std::string text = "programming";
    print_prefix(text, 7); // program
    print_prefix("hello", 3); // hel
}
\end{lstlisting}

\begin{lstlisting}[caption={Blad czasu zycia -- tylko ilustracja, nie uruchamiac.}]
std::string_view bad() {
    std::string local = "temporary";
    return local; // dangling view po zakonczeniu funkcji
}
\end{lstlisting}

Zmiana rozmiaru \code{std::string} może unieważnić wcześniejsze widoki.
Dodatkowo \code{string_view::substr} nie kopiuje znaków, natomiast
\code{string::substr} zwraca nowy \code{string}. W API akceptującym
tekst wejściowy dobrze jest dokumentować, kto posiada pamięć i jak
długo wywoływana funkcja może jej używać.

\subsection{UTF-8, Unicode i kodowanie}

\code{std::string} jest ciągiem jednostek typu \code{char}; sam z siebie
nie jest zbiorem znaków Unicode. W UTF-8 jeden punkt kodowy może
zajmować 1--4 bajty, a widoczny dla użytkownika znak (\emph{grapheme
cluster}) może składać się z kilku punktów kodowych. Dlatego
\code{text.size()} liczy bajty, a nie litery. \code{char8_t} (C++20)
reprezentuje jednostkę kodową UTF-8, lecz \code{u8"..."} ma typ
\code{const char8_t[N]}, nie \code{const char[N]}.

\begin{lstlisting}[caption={Rozmiar w bajtach i liczba punktow kodowych UTF-8.}]
#include <cstddef>
#include <iostream>
#include <string_view>

std::size_t codepoints_utf8(std::string_view bytes) {
    std::size_t n = 0;
    for (unsigned char c : bytes)
        if ((c & 0xc0u) != 0x80u) ++n;
    return n;
}

int main() {
    std::string_view text = "A\xC4\x85"; // A oraz 'a' z ogonkiem
    std::cout << text.size() << ' ' << codepoints_utf8(text) << '\n';
}
\end{lstlisting}

Funkcja z przykładu \emph{zakłada poprawne UTF-8} i liczy jedynie
jednostki niebędące bajtami kontynuacji: nie jest walidatorem kodowania.
Dla danych niezaufanych wymagane jest pełne sprawdzenie składni UTF-8,
a przy operacjach językowych --- biblioteka z obsługą normalizacji,
segmentacji i porównywania Unicode. \code{std::codecvt} nie jest zalecanym
nowym API (część jego komponentów została wycofana lub uznana za
przestarzałą).

\section{Parsowanie liczb przez \texttt{from\_chars}}
\label{sec:cpp19-charconv}

Funkcje \code{std::from_chars} i \code{std::to_chars} pracują na
przedziałach bufora, nie wymagają tworzenia strumienia i nie zależą
od lokalizacji \code{std::locale}. Wynik \code{from_chars} zawiera
wskaźnik na pierwszy nieprzetworzony znak oraz \code{std::errc}.
W szczególności udany odczyt prefiksu \code{123} z \code{123abc}
nie oznacza, że całe wejście jest poprawną liczbą.

\begin{lstlisting}[caption={Poprawny parser pelnej liczby calkowitej (C++17).}]
#include <charconv>
#include <optional>
#include <string_view>
#include <system_error>

std::optional<int> parse_int(std::string_view text) {
    int value{};
    auto result = std::from_chars(text.data(),
                                  text.data() + text.size(), value);
    if (result.ec != std::errc{} ||
        result.ptr != text.data() + text.size())
        return std::nullopt;
    return value;
}
\end{lstlisting}

Inaczej niż odczyt formatowany, \code{from_chars} nie pomija
początkowych białych znaków. Dla typów zmiennoprzecinkowych należy
uwzględnić dostępność odpowiedniej implementacji standardowej biblioteki
oraz wymagania dotyczące wartości nieskończonych, NaN i zaokrągleń.

\section{Parser pliku CSV: przykład kontraktu formatu}

CSV nie jest po prostu \code{split(',')}: poprawne pola mogą zawierać
przecinek i znaki cudzysłowu, a cytowane pola także nowe linie.
Najpierw należy określić dialekt, kodowanie, maksymalny rozmiar pola
i zachowanie wobec niepoprawnych danych. Poniższy prosty parser
obsługuje jeden rekord przechowywany już jako \code{string_view},
cudzysłowy i podwojone cudzysłowy, ale \emph{nie obsługuje wieloliniowych
rekordów} (ich odczyt wymaga osobnej warstwy).

\begin{lstlisting}[caption={Parser pojedynczego rekordu CSV (C++20).}]
#include <stdexcept>
#include <string>
#include <string_view>
#include <vector>
#include <utility>

std::vector<std::string> parse_csv_record(std::string_view s) {
    std::vector<std::string> fields;
    std::string field;
    bool quoted = false;
    bool just_closed = false;
    bool at_field_start = true;
    for (std::size_t i = 0; i < s.size(); ++i) {
        const char c = s[i];
        if (quoted) {
            if (c == '"') {
                if (i + 1 < s.size() && s[i + 1] == '"') {
                    field.push_back('"'); ++i;
                } else {
                    quoted = false;
                    just_closed = true;
                }
            } else {
                field.push_back(c);
            }
        } else if (c == ',' ) {
            fields.push_back(std::move(field));
            field.clear();
            at_field_start = true;
            just_closed = false;
        } else if (c == '"' && at_field_start) {
            quoted = true;
            at_field_start = false;
        } else {
            if (just_closed || c == '"')
                throw std::invalid_argument("invalid CSV quote");
            field.push_back(c);
            at_field_start = false;
        }
    }
    if (quoted) throw std::invalid_argument("unclosed quote");
    fields.push_back(std::move(field));
    return fields;
}
\end{lstlisting}

Istotne jest ograniczenie pamięci: parser przyjmujący dowolnie długi
rekord może przy dużym wejściu spowodować wyczerpanie zasobów.
Zachowanie \code{std::regex} również nie zastępuje pełnego parsera CSV;
wyrażenia regularne nie są ponadto narzędziem do ogólnej segmentacji Unicode.

\section{Wydajność i projektowanie interfejsów I/O}

Koszt operacji na nośniku i koszty systemowe często dominują nad
samymi obliczeniami. Nie należy jednak bez pomiaru utożsamiać
\code{mmap} z szybkością, a \code{fstream} z wolnym rozwiązaniem.
Przed optymalizacją warto mierzyć liczbę wywołań systemowych,
wielkość bufora, przepustowość, koszt parsowania i alokacji.

\begin{lstlisting}[caption={Wczytanie calego pliku binarnego z kontrola bledu (C++20).}]
#include <fstream>
#include <iterator>
#include <stdexcept>
#include <string>

std::string read_all(const std::string& path) {
    std::ifstream input(path, std::ios::binary);
    if (!input) throw std::runtime_error("open failed");
    std::string result((std::istreambuf_iterator<char>(input)),
                        std::istreambuf_iterator<char>{});
    if (input.bad()) throw std::runtime_error("read failed");
    return result;
}
\end{lstlisting}

Funkcja jest prosta i przydatna dla plików o kontrolowanym rozmiarze.
Jeśli plik może być wielogigabajtowy lub pochodzi od niezaufanego
nadawcy, lepiej przetwarzać go strumieniowo, z limitami i obsługą
częściowych rekordów. Dla operacji formatowania warto rozważyć
\code{std::format} (C++20) i \code{std::print} (C++23), ale przy
\code{std::print} należy sprawdzić wspierającą implementację.

\section{Konwersja liczb w obu kierunkach: \texttt{to\_chars}}

Dla protokołów tekstowych często potrzeba zamieniać liczby na tekst
w już przydzielonym buforze. \code{std::to_chars} nie tworzy
\code{std::string}, nie używa lokalizacji i informuje, gdzie skończył
pisać. Użytkownik odpowiada za wielkość bufora.

\begin{lstlisting}[caption={Formatowanie liczby calkowitej bez alokacji (C++17).}]
#include <array>
#include <charconv>
#include <stdexcept>
#include <string>
#include <system_error>

std::string format_integer(int value) {
    std::array<char, 32> buffer{};
    auto [ptr, ec] = std::to_chars(buffer.data(),
                                   buffer.data() + buffer.size(),
                                   value);
    if (ec != std::errc{})
        throw std::runtime_error("buffer too small");
    return std::string(buffer.data(), ptr);
}
\end{lstlisting}

Warto odróżnić brak alokacji w samej konwersji od alokacji przy
zbudowaniu wynikowego \code{std::string}. Można zamiast tego pisać
do bufora zarządzanego przez wywołującego; wtedy jego czas życia i
pojemność stają się elementem kontraktu.

\section{Małe protokoły binarne: odczyt kontrolowany}

Zapis wersjonowany nie musi oznaczać złożonej biblioteki. Jeżeli
format przewiduje czterobajtową liczbę little-endian, parsowanie
powinno sprawdzić, czy rzeczywiście otrzymaliśmy dokładnie tyle
bajtów. Należy jasno ustalić, czy nadmiarowe dane są dozwolone.

\begin{lstlisting}[caption={Odczyt 32-bitowej liczby little-endian (C++20).}]
#include <array>
#include <cstdint>
#include <istream>
#include <stdexcept>

std::uint32_t read_u32_le(std::istream& input) {
    std::array<unsigned char, 4> b{};
    for (auto& byte : b) {
        char c{};
        if (!input.get(c))
            throw std::runtime_error("truncated input");
        byte = static_cast<unsigned char>(c);
    }
    std::uint32_t result = 0;
    for (unsigned i = 0; i < 4; ++i)
        result |= std::uint32_t{b[i]} << (8 * i);
    return result;
}
\end{lstlisting}

W plikach o zmiennym formacie warto stosować pola: sygnaturę
(\emph{magic}), wersję, długość rekordu i jego zawartość. Jeśli
długość pochodzi z pliku, konieczny jest limit, kontrola przepełnienia
przy obliczaniu rozmiaru alokacji i obsługa przerwania odczytu.

\section{Wyjście wielowątkowe i \texttt{std::osyncstream}}

Wielu producentów logów może wywoływać \code{std::cout} jednocześnie.
Nawet jeśli biblioteka zapobiega wewnętrznemu wyścigowi danych,
poszczególne znaki wielu wiadomości mogą się przeplatać.
C++20 dostarcza \code{std::osyncstream}, który grupuje wypisywanie
i emituje zebrany fragment przy synchronizacji bufora lub destrukcji.

\begin{lstlisting}[caption={Niezachodzace na siebie linie logu (C++20).}]
#include <iostream>
#include <syncstream>
#include <thread>
#include <vector>

int main() {
    std::vector<std::jthread> workers;
    for (int id = 0; id < 4; ++id) {
        workers.emplace_back([id] {
            std::osyncstream(std::cout) << "worker " << id << '\n';
        });
    }
}
\end{lstlisting}

Kolejność wyświetlenia całych komunikatów nadal zależy od wykonania
wątków. \code{osyncstream} nie służy do synchronizacji obliczeń ani
do wymuszania kolejności zdarzeń w algorytmie.

\section{Szczegóły przenośności: ścieżki i kodowanie nazw plików}

Reprezentacja \code{filesystem::path} jest zależna od platformy.
Na Windows ścieżka może mieć natywną reprezentację opartą na szerokich
znakach, podczas gdy systemy POSIX posługują się sekwencjami bajtów.
Nie należy utożsamiać \code{path.string()} z konwersją Unicode
z identyczną semantyką na wszystkich platformach.

Również separator katalogów, wielkość liter, dozwolone nazwy,
obsługa ścieżek sieciowych i normalizacja są platformowe.
\code{lexically_normal()} upraszcza składnię ścieżki, lecz
nie wykonuje operacji I/O i nie rozstrzyga, dokąd prowadzą łącza
symboliczne. \code{canonical()} odnosi się do faktycznego systemu
plików, a zatem może zgłosić błąd.

\section{Projektowanie parserów i diagnostyka błędów}

Dobre API parsera zwraca nie tylko rezultat, ale -- gdy jest to
potrzebne -- także pozycję błędu. Format powinien określać,
czy dopuszcza nadmiarowe białe znaki, puste pola, separator na końcu
i dodatkowe dane. W warstwie parsera warto rozróżnić błędny zapis
liczby od jej przepełnienia oraz zbyt długi rekord od przedwczesnego
końca pliku. Takie rozróżnienie ułatwia generowanie precyzyjnych
komunikatów i testów regresyjnych.

W systemach obsługujących dane niezaufane kod przetwarzający tekst
powinien obowiązkowo mieć testy brzegowe: pusty plik, zerowy bajt,
wielobajtowy znak na granicy bufora, cytowanie CSV, rekord przekraczający
limit oraz zmianę pliku podczas odczytu. Dla parserów złożonych
pomocne są fuzzery i sanitizatory opisane w rozdziale 18.


\section{Podsumowanie: kiedy używać którego narzędzia?}
\begin{center}\small
\begin{tabularx}{\textwidth}{@{}p{0.32\textwidth}X@{}}
\toprule
Potrzeba & Narzędzie i uwagi \\
\midrule
Niewielkie pliki konfiguracyjne & \code{ifstream}, \code{getline}, oddzielny parser, walidacja. \\
Formatowanie komunikatów & \code{ostringstream} albo \code{format} przy dostępnej implementacji. \\
Duże przetwarzanie liczb tekstowych & \code{from_chars}, buforowanie, pomiar narzutów. \\
Ścieżki i metadane & \code{filesystem}, obsługa \code{error_code} i wyjątków. \\
API odczytujące tekst & \code{string_view}, pod warunkiem poprawnego czasu życia. \\
Trwałe dane binarne & Wersjonowany format; nie bezpośredni zapis struktur. \\
Unicode i operacje językowe & Walidacja UTF-8 i wyspecjalizowana biblioteka. \\
\bottomrule
\end{tabularx}
\end{center}

\nota{Trzy częste błędy to: \code{while(!eof())}, ignorowanie błędu
zapisu oraz zachowywanie \code{string_view} do zniszczonego bufora.
Z kolei w kodzie wrażliwym na bezpieczeństwo samo sprawdzenie, czy
ścieżka istnieje, nie zapewnia bezpiecznego otwarcia.}

\section*{Literatura i dokumentacja}
\addcontentsline{toc}{section}{Literatura i dokumentacja}
\begingroup\sloppy
\begin{itemize}
\item \url{https://en.cppreference.com/w/cpp/io} --- strumienie i ich stany.
\item \url{https://en.cppreference.com/w/cpp/filesystem} --- \code{filesystem}.
\item \url{https://en.cppreference.com/w/cpp/utility/from_chars} --- parsowanie liczb.
\item \url{https://en.cppreference.com/w/cpp/string/basic_string_view} --- widoki tekstowe.
\item \url{https://en.cppreference.com/w/cpp/utility/format/format} --- formatowanie.
\item The Unicode Consortium, \emph{The Unicode Standard}, \url{https://unicode.org/standard/standard.html}.
\end{itemize}
\endgroup
